%-----------------------------------------------------------------
% IDENTITAS MAHASISWA
%
% File ini adalah SATU-SATUNYA file yang perlu diubah setiap kali template
% ini dipakai. Semua halaman (sampul, pengesahan, pernyataan, kata pengantar)
% membaca data dari sini, sehingga tidak ada data pribadi yang tersebar di
% file lain.
%
% Cara pakai: setiap data memiliki dua macro.
%   \setnamadata{...}  diisi satu kali di file ini (setter)
%   \namadata          dibaca di seluruh naskah (getter)
% Contoh: \setfullname{Muhammad Ichsan} diisi di sini, sedangkan
% \fullname dipakai di tabel halaman pernyataan.
%
% Ganti nilai di dalam kurung kurawal { } dengan data Anda.
%-----------------------------------------------------------------


%-----------------------------------------------------------------
% 1. JUDUL TESIS
%-----------------------------------------------------------------

% Judul untuk halaman sampul (boleh lebih panjang/menggunakan \textit{})
\setjudulcover{[Judul Tesis dalam Bahasa Indonesia]}

% Judul untuk halaman pengesahan dan berita acara
\setjudul{[Judul Tesis dalam Bahasa Indonesia]}

% Judul dalam Bahasa Inggris
\setjudulinggris{[Thesis Title in English]}


%-----------------------------------------------------------------
% 2. IDENTITAS MAHASISWA
%-----------------------------------------------------------------

% Nama lengkap tanpa gelar, tidak boleh dicetak tebal
\setfullname{[Nama Lengkap Mahasiswa]}

% NIM (Nomor Induk Mahasiswa)
\setidnum{[0000000000000]}

% Tempat dan tanggal lahir (dipakai pada surat pernyataan bebas plagiasi)
\settempatlahir{[Kota/Kabupaten]}
\settanggallahir{[dd bulan yyyy]}

% Status akademik (biasanya: Mahasiswa)
\setstatusmahasiswa{Mahasiswa}


%-----------------------------------------------------------------
% 3. PROGRAM STUDI DAN INSTANSI
%-----------------------------------------------------------------

% Nama program studi, contoh: Magister Kecerdasan Buatan
\setprogramstudi{Magister Kecerdasan Buatan}

% Nama departemen tempat program studi berada
\setdept{Informatika}

% Fakultas, universitas, dan kota pencetakan
\setfaculty{Fakultas Matematika dan Ilmu Pengetahuan Alam}

% Nama singkat fakultas untuk halaman pengesahan
\setfakultassingkat{FMIPA}
\setuniversity{Universitas Syiah Kuala}
\setcity{Banda Aceh}

% Gelar yang diperoleh, contoh: Magister Komputer (M.Kom.)
% SESUAIKAN dengan gelar resmi program studi Anda.
\setjenjang{Magister}

\setgelar{Magister Komputer}

% Jenis dokumen: TESIS
\setdokumen{TESIS}

% Kalimat sebelum gelar di sampul.
% Contoh untuk program magister: "Diajukan untuk memenuhi sebagian
% persyaratan memperoleh gelar"
\setsyarattext{Diajukan untuk melengkapi tugas-tugas dan\\ memenuhi syarat-syarat guna memperoleh gelar}

% Tahun penulisan naskah tesis
\setyearsubmit{[Bulan], [Tahun]}


%-----------------------------------------------------------------
% 4. DOSEN PEMBIMBING
%-----------------------------------------------------------------

% Pembimbing I
\setfirstsupervisor{[Nama Dosen Pembimbing I]}
\setfirstnip{[NIP Pembimbing I]}

% Pembimbing II
\setsecondsupervisor{[Nama Dosen Pembimbing II]}
\setsecondnip{[NIP Pembimbing II]}


%-----------------------------------------------------------------
% 5. PIMPINAN FAKULTAS DAN PROGRAM STUDI
%
% Isi dengan nama pejabat yang berlaku pada saat naskah
% dicetak. Cek ke bagian/program studi bila sudah berubah.
%-----------------------------------------------------------------

% Dekan fakultas
\setdekan{[Nama Dekan]}
\setdekannip{[NIP Dekan]}

% Ketua departemen (jika ditampilkan pada halaman pengesahan)
\setkajur{[Nama Ketua Departemen]}
\setkajurnip{[NIP Ketua Departemen]}

% Ketua program studi
\setkaprodi{[Nama Ketua Program Studi]}
\setkaprodinip{[NIP Ketua Program Studi]}

% Koordinator tesis (tanda tangan pada surat pernyataan)
\setkoordinatortesis{[Nama Koordinator Tesis]}
\setkoordinatortesisnip{[NIP Koordinator Tesis]}


%-----------------------------------------------------------------
% 6. TANGGAL
%-----------------------------------------------------------------

% Tanggal ujian tesis, ditulis lengkap dengan nama hari.
% Contoh: Senin, 17 Februari 2025
\setapprovaldate{[Hari], [Tanggal] [Bulan] [Tahun]}

% Tanggal penanda tanganan pada kata pengantar, surat pernyataan, dan
% kalimat "Tesis disahkan pada hari ...".
% Contoh: 14 Februari 2025
\setdisahkantanggal{[Tanggal] [Bulan] [Tahun]}

% Nama hari untuk kalimat "Tesis disahkan pada hari ..."
% Contoh: Senin
\setdisahkantanggalhari{[Hari]}


%-----------------------------------------------------------------
% 7. KATA KUNCI ABSTRAK (opsional, dipakai sebagai pengingat)
%-----------------------------------------------------------------
\setkeywordindo{[Kata kunci 1, Kata kunci 2, Kata kunci 3]}
\setkeywordeng{[Keyword 1, Keyword 2, Keyword 3]}


%-----------------------------------------------------------------
% End of identitas.tex
%-----------------------------------------------------------------